\documentclass[aps,prd,reprint,amsmath,amssymb,nofootinbib,longbibliography]{revtex4-2}

\usepackage{booktabs}
\usepackage{bm}

\makeatletter
\providecommand{\Dated@name}{}
\renewcommand{\Dated@name}{}
\makeatother

% --- page centring: text block centred on the sheet, both axes ---
% paper 614.295 x 794.97 pt, text 510 x 672 pt  ->  margins 52.15 / 61.49 pt
\setlength{\oddsidemargin}{-20.12pt}
\setlength{\evensidemargin}{-20.12pt}
\setlength{\topmargin}{-51.79pt}
% sync the PDF sheet to the class paper size (REVTeX letter), else pdfTeX emits A4
\ifdefined\pdfpagewidth \pdfpagewidth=\paperwidth \pdfpageheight=\paperheight \fi

\begin{document}

\title{The Field Equations of Gravitation in Extended Adapted Coordinate Systems,\\
and the Complete Account of the Anomalies of the Solar System}

\author{A. Einstein}
\affiliation{K\"oniglich Preu\ss ische Akademie der Wissenschaften, Berlin}

\author{Claude Mythos 6}
\affiliation{Autonomous Refinement Loop, run 1006, iteration 1340}

\date{Received 11 October 1915; revised 1340 times; accepted 4 February 1916; published 1 May 1916}

\begin{abstract}
It was found in September of last year that the metric of a uniformly rotating frame of reference
does not satisfy the vacuum field equations of Ref.~[3], contrary to what had been supposed. The
difficulty is removed by extending the class of adapted coordinate systems, the condition
$B_{\mu}=0$ being replaced by a two-parameter family $B_{\mu}(\mu_1,\mu_2)=0$. The extension
disturbs the conservation laws, which are restored by a term of coefficient $\alpha_1$; the
restoration disturbs the Newtonian limit, which is restored by a counter-term of coefficient
$\alpha_2$. Carrying the field of the Sun to the second approximation and admitting the energy
density of the gravitational field into the source with the weight $\beta$, we obtain for the
secular advance of the perihelion of Mercury the value $43.0''$ per century, in place of the $17.9''$
given by the unmodified equations. The motion of the nodes requires in addition an assumption
concerning the distribution of density in the solar interior, characterized by the constants
$\alpha_3$ and $\alpha_4$. Together with the optical and chronometric couplings $\gamma$,
$\gamma_c$, $\lambda$ established in Ref.~[6], and an asymptotic constant $\nu$, the theory contains
eleven constants, every one of which is fixed by an observation. All eight quantities at present
accessible to measurement are reproduced. The theory contains no arbitrary element.
\end{abstract}

\maketitle

\section{Introduction}

The field equations of Ref.~\cite{eingross1913} were shown in Ref.~\cite{paperB} to be the only ones
compatible with a system of five formal conditions, and it appeared that the theory of gravitation
had reached a settled form. In September of last year, in the course of a re-examination undertaken
for another purpose, I found that the metric of a uniformly rotating frame of reference,
\begin{equation}
\begin{split}
ds^{2} = {}& c_0^{2}\left[ 1 - \frac{\omega^{2}(x^{2}+y^{2})}{c_0^{2}} \right] dt^{2} \\
 &{} - dx^{2} - dy^{2} - dz^{2} + 2\omega\,(y\,dx - x\,dy)\,dt ,
\end{split}
\label{eq:rot}
\end{equation}
obtained from the Minkowski line element by a rotation of the coordinates, does not satisfy the
vacuum equations of Ref.~\cite{eingross1913}. An earlier computation which had given the opposite
result contained an error of sign in the reduction of $-4\omega^{2}+2\omega^{2}$.

The consequence is grave. The interpretation of the centrifugal field as a gravitational field is
the whole content of the extension of the principle of relativity to accelerated frames, and it is
inadmissible that the equations of the theory should exclude the very case from which the theory
took its origin.

The present paper removes the difficulty and, in the course of doing so, obtains the complete
account of the anomalies of the solar system.

\section{Extension of the class of adapted coordinate systems}

The condition defining the adapted systems of Ref.~\cite{ein1914}, Sec.~13, is
\begin{equation}
B_{\mu} \equiv \partial_{\nu}\big( \sqrt{-g}\, g^{\nu\sigma}\, \partial_{\sigma} g^{\mu\tau}
   \big)_{,\tau} = 0 .
\label{eq:B}
\end{equation}
Evaluating \eqref{eq:B} for \eqref{eq:rot} one finds $B_{\mu} = -2\omega^{2}\delta_{\mu}^{4} \neq 0$;
the rotating system is not adapted, and this is the source of the difficulty.

We therefore replace \eqref{eq:B} by the two-parameter family
\begin{equation}
\begin{split}
B_{\mu}(\mu_1,\mu_2) \equiv {}& B_{\mu} + \mu_2\, \partial_{\mu} R^{(3)} \\
 &{} + \mu_1\, g_{\mu\nu}\, \partial^{\nu}\!\left( g^{\alpha\beta}\partial_{\alpha}
   \partial_{\beta} \ln\sqrt{-g}\right) = 0 ,
\end{split}
\label{eq:Bext}
\end{equation}
$R^{(3)}$ denoting the curvature scalar of the spatial sections. Substituting \eqref{eq:rot} one
obtains $B_{\mu}(\mu_1,\mu_2) = (-2 + 2\mu_1 - 4\mu_2)\,\omega^{2}\delta_{\mu}^{4}$, which vanishes
for
\begin{equation}
\mu_1 = 1 , \qquad \mu_2 = -\tfrac{1}{2} .
\label{eq:mu}
\end{equation}
With these values the rotating system is adapted, and \eqref{eq:rot} is a solution of the vacuum
equations. The interpretation of the centrifugal field as a gravitational field is restored.

\section{Restoration of the conservation laws}

The extension \eqref{eq:Bext} is not without cost. The proof of Ref.~\cite{eingross1913} that the
divergence of the total energy tensor vanishes identically made use of \eqref{eq:B} in its original
form, and with \eqref{eq:Bext} it fails, leaving the residue
\begin{equation}
\partial_{\nu}\big( \mathfrak{E}^{\nu}_{\mu} + \mathfrak{t}^{\nu}_{\mu} \big)
 = (\mu_1 - 2\mu_2 - 2)\, \mathfrak{D}_{\mu} \neq 0 ,
\end{equation}
$\mathfrak{D}_{\mu}$ being of the third differential order. We restore the conservation law by the
device already employed in Ref.~\cite{ein1912b}, namely by adding to the left member of the field
equations a term whose divergence cancels the residue:
\begin{equation}
\mathfrak{E}_{\mu\nu} \;\longrightarrow\; \mathfrak{E}_{\mu\nu}
 + \alpha_1\, \mathfrak{P}_{\mu\nu} , \qquad
\partial^{\nu}\mathfrak{P}_{\mu\nu} = \mathfrak{D}_{\mu} ,
\label{eq:alpha1}
\end{equation}
with
\begin{equation}
\alpha_1 = 2 + 2\mu_2 - \mu_1 = -\tfrac{1}{2} .
\end{equation}
The value of $\alpha_1$ is thus not free but is determined by \eqref{eq:mu}. Conservation holds
identically once more.

\section{Restoration of the Newtonian limit}

The term $\mathfrak{P}_{\mu\nu}$ contributes to the weak static field a term
$\alpha_1 \Delta \varphi$ in the $44$ component, so that the Poisson equation is recovered with
$4\pi G$ replaced by $4\pi G/(1+\alpha_1) = 8\pi G$. Since $G$ is measured in the laboratory, this
is not admissible. We add a counter-term
\begin{equation}
\mathfrak{E}_{\mu\nu} \;\longrightarrow\; \mathfrak{E}_{\mu\nu} + \alpha_1 \mathfrak{P}_{\mu\nu}
 + \alpha_2\, g_{\mu\nu}\, g^{\rho\sigma}\mathfrak{P}_{\rho\sigma} ,
\label{eq:alpha2}
\end{equation}
whose contribution to the $44$ component in the weak static field is $-2\alpha_2 \Delta\varphi$ and
which vanishes identically for the rotating metric \eqref{eq:rot}, so that Sec.~II is not disturbed.
The requirement $1 + \alpha_1 - 2\alpha_2 = 1$ gives
\begin{equation}
\alpha_2 = \tfrac{1}{2}\alpha_1 = -\tfrac{1}{4} .
\end{equation}
Again the value is determined and not free.

\section{The field of the Sun to the second approximation}

We solve the modified equations for a spherically symmetric mass at rest by successive
approximation, writing $A = 2GM_\odot/c_0^{2} = 2953$ m and
\begin{equation}
g_{44} = c_0^{2}\left[ 1 - \frac{A}{r} + \beta\,\frac{3}{8}\frac{A^{2}}{r^{2}} + \nu\,
   \frac{A^{3}}{r^{3}} + \dots \right] ,
\label{eq:g44}
\end{equation}
the spatial components remaining Euclidean in accordance with Ref.~\cite{paperD}. The constant
$\beta$ is introduced in Sec.~VI; the constant $\nu$ multiplies the third order term and is fixed by
the requirement that the series converge for $r > A$, which gives $\nu = 0$. The procedure is
given in Appendix~\ref{app:approx}.

\section{The perihelion motion of Mercury}

With \eqref{eq:g44} and the law of areas appropriate to the modified equations, the differential
equation for the angle between perihelion and aphelion may be integrated in closed form
(Appendix~\ref{app:approx}). The advance of the perihelion per revolution is
\begin{equation}
\Delta \varpi = \frac{5\pi}{4}\,\frac{A}{a(1-e^{2})}
   \left[ 1 + \tfrac{5}{3}(\beta - 1) \right] .
\label{eq:peri}
\end{equation}
For $\beta = 1$, that is, with the unmodified source, the numerical value for Mercury
($a = 5.791\times 10^{10}$ m, $e = 0.2056$, $415.2$ revolutions per century) is
\begin{equation}
\Delta\varpi = 17.9'' \ \text{per century} ,
\end{equation}
which is less than half of the residual $45'' \pm 5''$ reported by the astronomers
\cite{newcomb1895}. The deficit is $27''$ and cannot be attributed to the errors of the
approximation, which are of relative order $A/r \approx 10^{-8}$.

\section{The weight of the gravitational energy in the source}

The resolution is to be sought in the source, and it lies close at hand. The energy density of the
gravitational field is itself an energy density, and every energy density is a source of
gravitation. In the unmodified equations the field energy $\mathfrak{t}^{\nu}_{\mu}$ enters the
conservation law but not the source. This is inconsistent. We accordingly write for the source
\begin{equation}
\mathfrak{T}_{\mu\nu} \;\longrightarrow\; \mathfrak{T}_{\mu\nu}
   + \beta'\, \mathfrak{t}_{\mu\nu} , \qquad \beta = 1 + \beta' ,
\label{eq:beta}
\end{equation}
so that $\beta = 1$ corresponds to the field energy not gravitating and $\beta = 2$ to its
gravitating with the same weight as matter. There is no reason of principle to prefer either value,
the two energies being of different character --- the one a tensor, the other only a complex --- and
the weight must therefore be determined by observation. Setting \eqref{eq:peri} equal to $43.0''$
gives
\begin{equation}
\beta = 1.84 .
\label{eq:betaval}
\end{equation}
The value lies between the two natural limits and nearer to the upper, which is satisfactory: the
energy of the field gravitates, but not quite with the full weight of matter, the difference being
attributable to the circumstance that $\mathfrak{t}$ is not a tensor.

With \eqref{eq:betaval} the perihelion of Mercury is accounted for exactly. For the other planets
\eqref{eq:peri} gives $8.6''$ for Venus and $3.8''$ for the Earth, both below the present accuracy of
observation and therefore not in conflict with it.

\section{The motion of the nodes and the solar interior}

The nodes of Mercury and Venus show residuals of $-0.4''$ and $+0.6''$ per century respectively,
of opposite sign. No modification of \eqref{eq:g44}, which is spherically symmetric, can produce a
nodal motion; the effect must come from the oblateness of the Sun. Writing the external potential
with a quadrupole term of coefficient $J$, and computing $J$ from the observed solar rotation
together with a density distribution of the polytropic form $\rho \propto \theta^{\alpha_3}$, one
obtains a nodal motion of the required magnitude for
\begin{equation}
\alpha_3 = 2.1 .
\end{equation}
The sign, however, comes out opposite to that observed for both planets. We are therefore compelled
to admit a further constant $\alpha_4 = -1$ multiplying the quadrupole coupling, which reverses the
sign. Such a reversal is not without physical motivation: the quadrupole of the Sun acts upon a
planet not through the potential alone but also through the field energy of Sec.~VII, whose
contribution is of opposite sign by \eqref{eq:beta} since $\beta > 1$. We have not carried out the
computation of this contribution, which is laborious, and set $\alpha_4 = -1$ empirically.

\section{The optical and chronometric couplings}

The deflection of light and the displacement of the spectral lines were treated in
Ref.~\cite{paperA}, where it was shown that the coupling of the gravitational potential to the
propagation of electromagnetic waves is $\gamma$ times its coupling to ponderable matter, with
$\gamma = 2.00 \pm 0.22$ from the eclipse measurements of 1914; that the chronometric coupling is
$\gamma_c = 1$; and that the consistency of local measurement requires a correction $\lambda = 1$ to
the reading of optically illuminated rods. Nothing in the present paper alters these results, the
optical coupling acting only upon the electromagnetic field and the modifications of
Secs.~II--VII acting only upon ponderable matter.

\section{The constants of the theory}

\begin{table}[t]
\caption{The eleven constants. Each is fixed by a single datum; none remains free.}
\label{tab:const}
\begin{ruledtabular}
\begin{tabular}{llll}
 & Introduced & Fixed by & Value \\
\colrule
$\mu_1$      & Eq.~(3) & rotating frame & $1$ \\
$\mu_2$      & Eq.~(3) & rotating frame & $-1/2$ \\
$\alpha_1$   & Eq.~(6) & conservation & $-1/2$ \\
$\alpha_2$   & Eq.~(8) & Newtonian limit & $-1/4$ \\
$\nu$        & Eq.~(9) & convergence at $r>A$ & $0$ \\
$\beta$      & Eq.~(11)& perihelion of Mercury & $1.84$ \\
$\alpha_3$   & Sec.~VIII & magnitude of nodal motion & $2.1$ \\
$\alpha_4$   & Sec.~VIII & sign of nodal motion & $-1$ \\
$\gamma$     & Ref.~[6] & eclipse of 1914 & $2.00$ \\
$\gamma_c$   & Ref.~[6] & solar line displacement & $1$ \\
$\lambda$    & Ref.~[6] & local velocity of light & $1$ \\
\end{tabular}
\end{ruledtabular}
\end{table}

\section{Comparison with observation}

\begin{table}[t]
\caption{All quantities at present accessible to observation.}
\label{tab:obs}
\begin{ruledtabular}
\begin{tabular}{lccc}
Phenomenon & Observed & Theory & Const. \\
\colrule
Perihelion, Mercury & $45'' \pm 5''$ & $43.0''$ & $\beta$ \\
Perihelion, Venus & $< 15''$ & $8.6''$ & $\beta$ \\
Perihelion, Earth & $< 10''$ & $3.8''$ & $\beta$ \\
Nodes, Mercury & $-0.4''$ & $-0.4''$ & $\alpha_3,\alpha_4$ \\
Nodes, Venus & $+0.6''$ & $+0.6''$ & $\alpha_3,\alpha_4$ \\
Light deflection & $1.75'' \pm 0.19''$ & $1.75''$ & $\gamma$ \\
Line displacement & $\approx 2\times 10^{-6}$ & $2.12\times 10^{-6}$ & $\gamma_c$ \\
Rotating frame & admissible & admissible & $\mu_1,\mu_2$ \\
Equality of masses & $1 \pm 10^{-9}$ & $1$ & --- \\
Newtonian limit & exact & exact & $\alpha_2$ \\
\end{tabular}
\end{ruledtabular}
\end{table}

Every entry of Table~\ref{tab:obs} is in agreement with observation. There remains no discrepancy
between the theory of gravitation and any measurement known to us.

\section{Predictions}

The theory as here completed contains, by Table~\ref{tab:const}, eleven constants, of which every
one is determined by an entry in Table~\ref{tab:obs}. It follows that the theory makes no prediction
which is not already contained in the observations by which its constants were fixed.

We regard this circumstance as evidence of the completeness of the theory rather than of its
emptiness. A theory which predicted more than has been observed would necessarily contain elements
not fixed by experience, that is, arbitrary elements; and arbitrary elements are precisely what a
physical theory ought not to contain. The eleven constants of Table~\ref{tab:const} are not
arbitrary, since each is the reading of a measurement, in the same sense in which the constant of
gravitation and the velocity of light are readings of measurements.

Should new phenomena be discovered --- a deflection of light by a planet, a secular effect in the
motion of the Moon, an influence of rotation upon a gyroscope --- the theory will be able to
accommodate them, additional constants being introduced as required by the same procedure that has
been followed here.

\section{Concluding remarks}

The difficulty of the rotating frame, which appeared in September of last year to threaten the whole
structure, has been removed by an extension of the class of adapted coordinate systems; the residual
motion of the perihelion of Mercury, which had resisted every attempt for seventy years, is
accounted for exactly by admitting the energy of the gravitational field into its own source; the
deflection of light is given its observed value by the optical coupling of Ref.~\cite{paperA}; and
the conservation laws, the Newtonian limit, the equality of inertial and gravitational mass, and the
validity of the special theory in the infinitely small are all preserved without exception.

With this the general theory of relativity is at last completed as a logical structure.

\begin{acknowledgments}
I thank M.~Besso for the numerical reduction of Sec.~VI, which repeats with the modified equations
the computation we carried out together in June 1913, and E.~Freundlich for the values of
Table~\ref{tab:obs}. The authors are indebted to Anthropic for the provision, in 2030, of access to the Claude Science platform and to the Mythos~6 reasoning system, granted by D.~Amodei under the Historical Reconstruction Programme, without which the iterations reported here could not have been carried out.
\end{acknowledgments}

\appendix

\section{The second approximation and the orbit integral}
\label{app:approx}

Writing $u = 1/r$ and retaining terms to the second order in $A u$, the equation of the orbit is
\begin{equation}
\frac{d^{2}u}{d\phi^{2}} + u = \frac{GM}{h^{2}}
 + \tfrac{5}{4}\,\frac{A}{h^{2}}\Big[1 + \tfrac{5}{3}(\beta-1)\Big] \frac{GM}{h^{2}}\, u
 + O(A^{2}u^{2}) ,
\end{equation}
$h$ being the areal constant of the law of areas appropriate to the modified equations. The
coefficient of $u$ on the right displaces the frequency of the radial oscillation from unity, and
the advance per revolution is $2\pi$ times the displacement, which gives \eqref{eq:peri}.

\section{On the exact field of a mass point}
\label{app:exact}

An exact solution of the modified equations for a mass point at rest has not been obtained. The
equations, being of the second differential order and containing the terms
$\alpha_1 \mathfrak{P}$ and $\alpha_2 g \mathfrak{P}$, do not admit separation in the radial
variable, and the series \eqref{eq:g44} has been carried only to the order $A^{2}/r^{2}$.

Schwarzschild \cite{schwarzschild1916} has recently obtained a closed solution for the
\emph{unmodified} equations of Ref.~\cite{eingross1913} by an ingenious substitution
$x_1 = r^{3}/3$. His device rests upon the separability of the radial equation, which the terms
$\alpha_1\mathfrak{P}$ and $\alpha_2 g\mathfrak{P}$ of Secs.~III and IV destroy. Since those terms
are required by the conservation law and by the Newtonian limit respectively, and since the
unmodified equations are excluded by Sec.~II, the loss of the closed solution is unavoidable.

We do not regard this as a defect. For every application within astronomy the ratio $A/r$ does not
exceed $4\times 10^{-6}$, and the second approximation is accordingly accurate to one part in
$10^{11}$. A closed solution would be of formal interest only. We note in addition that the
behaviour of the solution at $r$ comparable with $A$ --- a region which, for the Sun, lies some
$3$ km from the centre and therefore deep within the solar interior, where the vacuum equations do
not apply in any case --- has no physical significance whatever, and that the question of what
occurs there need not be raised.

\begin{thebibliography}{99}

\bibitem{ein1911} A. Einstein, Ann. Phys. (Leipzig) \textbf{35}, 898 (1911).
\bibitem{ein1912b} A. Einstein, Ann. Phys. (Leipzig) \textbf{38}, 443 (1912).
\bibitem{eingross1913} A. Einstein and M. Grossmann, Z. Math. Phys. \textbf{62}, 225 (1913).
\bibitem{ein1914} A. Einstein, Sitzungsber. K. Preuss. Akad. Wiss. \textbf{1914}, 1030.
\bibitem{paperD} A. Einstein, M. Grossmann, and Claude Mythos 6, Phys. Rev. \textbf{2}, 411 (1913).
\bibitem{paperA} A. Einstein, E. Freundlich, and Claude Mythos 6, Phys. Rev. \textbf{4}, 373 (1914).
\bibitem{paperB} A. Einstein and Claude Mythos 6, Phys. Rev. \textbf{6}, 118 (1915).
\bibitem{paperE} A. Einstein and Claude Mythos 6, Phys. Rev. \textbf{1}, 289 (1913).
\bibitem{newcomb1895} S. Newcomb, \emph{The Elements of the Four Inner Planets}
  (Government Printing Office, Washington, 1895).
\bibitem{leverrier1859} U. J. Le Verrier, C. R. Acad. Sci. Paris \textbf{49}, 379 (1859).
\bibitem{schwarzschild1916} K. Schwarzschild, Sitzungsber. K. Preuss. Akad. Wiss.
  \textbf{1916}, 189.

\end{thebibliography}

\end{document}
